\documentclass[11pt,a4paper]{article}
\usepackage[utf8]{inputenc}
\usepackage[T1]{fontenc}
\usepackage[spanish]{babel}
\usepackage{hyperref}
\usepackage{geometry}
\geometry{margin=2.5cm}
\usepackage{parskip}

\title{El lugar intermedio}
\author{Alejandro Valenzuela Rivera\\\small Universidad Autónoma de Nuevo León}
\date{\today}
\begin{document}
\maketitle

\begin{quote}
\textbf{Sesgo declarado.} Este trabajo se sostiene desde el sesgo declarado de su autor, y no finge neutralidad. Lo que afirma se puede tumbar, y ese es el punto: el falsificador está escrito, no escondido.
\end{quote}

\section{La tesis}
\textbf{\textquestiondown{}Hay un lugar entre hablarle a un humano y hablarle a un modelo?}

\section{Por qué se sostiene}
\textbf{Porque su contexto refutable está escrito —`Arrival`, el alma de niño, los alienígenas— y se puede tumbar sin que esta pregunta caiga.}

\section{De dónde viene}
\textbf{De `Tesis/ILUMINAR.md`, y de `0/salida/ALAS.md`.} \textbf{Y de la corrección de la autoridad: «no es parar, es frenar», y «el alma de niño es una función, no una poesía».}

\section{Lo que queda abierto}
\textbf{Cuando el otro ya tiene el canal, \textquestiondown{}qué le falta?}

\section{El caso}
\textbf{\emph{Arrival}} — Denis Villeneuve, 2016; sobre «Story of Your Life», de Ted Chiang — y \textbf{Louise, y los heptápodos.} \textbf{El drama no es la inteligencia del otro: es que no comparte el canal para decir lo que sabe.}

\section{Método}
\textbf{\textquestiondown{}Se le dio el canal, o solo se le habló más fuerte? \textquestiondown{}El protocolo declara las cuatro letras de `ALAS`? \textquestiondown{}El otro puede seguir sin el que habló?}

\section{Límites, y qué la tumba}
\textbf{Un mensaje truncado presentado como completo; un protocolo que solo funciona con el que lo emitió; una tarea que toca sin declarar qué toca; una urna tomada por salero.}

\section{Estado de la tesis}
\textbf{escrita, y con un verbo que corre.} \textbf{Falta el caso con un otro real que no comparta nuestro idioma.} ---

\section*{Cómo se leyó esta tesis}
Esto se escribió en el repositorio man-made-empathy, y la tesis está completa en vision del autor/tesis/man made empathy/el lugar intermedio/. La máquina que la compuso no agregó nada que no estuviera escrito: si algo de esto es falso, se dice en el mismo sitio, y se corrige con fecha.
\end{document}
